\documentclass[11pt,a4paper]{article}

% ---------------------------------------------------------
% PACKAGES
% ---------------------------------------------------------
\usepackage[utf8]{inputenc}
\usepackage[T1]{fontenc}
\usepackage{lmodern}
\usepackage{amsmath, amssymb, amsfonts}
\usepackage{bm}
\usepackage{geometry}
\usepackage{graphicx}
\usepackage{hyperref}
\usepackage{authblk}
\usepackage{cite}
\usepackage{physics}
\usepackage{mathtools}
\usepackage{tensor}
\usepackage{enumitem}

\geometry{margin=2.4cm}

% ---------------------------------------------------------
% TITLE
% ---------------------------------------------------------
\title{\textbf{RDCC--Euclid Interface Paper V2.0:\\
A Survey-Ready Framework for Testing Relaxation-Driven Cosmology}}

\author[1]{Michael Lehmann}
\affil[1]{\small Independent Researcher, Relaxation-Driven Cyclic Cosmology (RDCC) Framework\\
\texttt{mi.lehmann@gmx.de}}

\date{\small June 2026}

\begin{document}
\maketitle

% ========================================================
% ABSTRACT
% ========================================================

\begin{abstract}
The Relaxation-Driven Cyclic Cosmology (RDCC) modifies the gravitational potential
through a scale- and redshift-dependent relaxation kernel $\mathcal{R}(k,z;\alpha)$.
These modifications induce geometric, morphological, and topological signatures in
weak-lensing observables that are directly accessible to Euclid's survey pipeline.
This paper provides a survey-ready, Euclid-focused overview of RDCC. We present:
(i) the theoretical structure of the relaxation kernel, (ii) the mapping from RDCC
modifications to Euclid observables, (iii) Euclid-level forecasts for the relaxation
parameter $\alpha$, (iv) the degeneracy structure and multi-probe synergies, and
(v) a practical integration pathway for incorporating RDCC into Euclid's forward
modeling, summary-statistics, and likelihood frameworks. The goal of this document
is to provide a concise, technically precise reference for Euclid collaborators and
theorists working on survey-level tests of relaxation-driven cosmology.
\end{abstract}

% ========================================================
% 1 INTRODUCTION
% ========================================================

\section{Introduction}

The Euclid mission enables high-precision tests of gravitational physics through
weak-lensing and galaxy-clustering measurements across a wide range of angular
scales and redshifts. The Relaxation-Driven Cyclic Cosmology (RDCC) provides a
minimal extension of $\Lambda$CDM in which the gravitational potential responds
to large-scale structure through a relaxation mechanism governed by a single
parameter $\alpha$. This mechanism introduces scale-dependent deviations in the
lensing potential that propagate into geometric and topological features of the
convergence field.

Euclid's tomographic binning, angular resolution, and survey depth make it
particularly sensitive to RDCC-induced distortions. These distortions appear in
power spectra, peak and void statistics, Minkowski functionals, persistent
homology, and transport-based topology. The combination of these probes allows
Euclid to constrain $\alpha$ at the percent level.

This paper provides a survey-ready interface between RDCC theory and Euclid
analysis pipelines. Section~2 summarizes the theoretical structure of RDCC.
Section~3 describes RDCC signatures in Euclid observables. Section~4 presents
Euclid-level forecasts. Section~5 analyzes degeneracies and multi-probe synergies.
Section~6 outlines integration pathways into Euclid pipelines. Section~7 provides
a concise conclusion.

% ========================================================
% NOMENCLATURE
% ========================================================

\section*{Nomenclature}

\begin{tabular}{ll}
$\alpha$ & RDCC relaxation parameter; controls scale-dependent modifications of the potential \\[4pt]
$\Phi$ & Gravitational potential \\[4pt]
$\Phi_{\Lambda\mathrm{CDM}}$ & Standard potential in the $\Lambda$CDM model \\[4pt]
$\Phi_{\mathrm{RDCC}}$ & Modified potential in the RDCC model \\[4pt]
$\mathcal{R}(k,z;\alpha)$ & RDCC relaxation kernel; scale- and redshift-dependent modification \\[4pt]
$k$ & Wavenumber in Fourier space \\[4pt]
$z$ & Redshift \\[4pt]
$\kappa$ & Weak-lensing convergence field \\[4pt]
$\gamma$ & Shear field \\[4pt]
$P_{\kappa}(\ell)$ & Convergence power spectrum as a function of multipole $\ell$ \\[4pt]
$\mathrm{MF}_{i}$ & Minkowski functionals (area, boundary length, Euler characteristic) \\[4pt]
$\mathcal{D}$ & Persistence diagram in persistent homology \\[4pt]
$W_{2}$ & Quadratic optimal-transport distance between persistence diagrams \\[4pt]
$\sigma(\alpha)$ & Forecasted uncertainty on the RDCC parameter $\alpha$ \\[4pt]
$\Delta \log Z$ & Evidence contrast between RDCC and $\Lambda$CDM \\[4pt]
$v_{k}$ & Fisher channel (eigenvector of the information-tensor decomposition) \\[4pt]
$\lambda_{k}$ & Eigenvalue of the $k$-th Fisher channel (information contribution) \\[4pt]
\end{tabular}

% ========================================================
% 2 RDCC: THEORETICAL OVERVIEW
% ========================================================

\section{RDCC: Theoretical Overview}

RDCC modifies the gravitational potential through a relaxation kernel
\begin{equation}
    \Phi_{\mathrm{RDCC}}(k,z)
    =
    \Phi_{\Lambda\mathrm{CDM}}(k,z)\,
    \mathcal{R}(k,z;\alpha),
\end{equation}
where $\mathcal{R}$ introduces scale-dependent deviations that become relevant
at $k \gtrsim 0.1\,h\,\mathrm{Mpc}^{-1}$. These deviations propagate into the
lensing potential and affect a wide range of weak-lensing observables, including
the convergence field $\kappa$, shear correlations, peak and void statistics,
Minkowski functionals, persistent homology, and transport-based topology.

RDCC is fully compatible with Euclid's tomographic binning and angular
resolution. Its signatures are strongest at intermediate multipoles accessible
to Euclid's VIS instrument.

% ========================================================
% 3 RDCC SIGNATURES IN EUCLID OBSERVABLES
% ========================================================

\section{RDCC Signatures in Euclid Observables}

RDCC induces characteristic distortions in Euclid weak-lensing data. These
distortions arise from the scale-dependent modification of the lensing potential
and propagate into geometric, morphological, and topological observables.

\subsection{Power Spectra}

The convergence power spectrum $P_{\kappa}(\ell)$ exhibits moderate-scale
suppression or enhancement depending on the value of $\alpha$. These effects are
partially degenerate with $\Lambda$CDM parameters but remain distinguishable
through tomographic and multi-probe combinations.

\subsection{Peak and Void Statistics}

RDCC shifts peak heights and void depths in the convergence field. These
statistics provide strong sensitivity to relaxation-driven potential changes and
are robust against shape noise at Euclid depth.

\subsection{Minkowski Functionals}

The Minkowski functionals $\mathrm{MF}_{i}$ capture morphological distortions in
area, boundary length, and Euler characteristic. RDCC induces characteristic
departures from the Gaussian predictions of $\Lambda$CDM.

\subsection{Persistent Homology}

RDCC modifies the birth–death structure of topological features in the
convergence field. Persistent homology provides a noise-resistant probe of these
modifications and is highly sensitive to the relaxation kernel.

\subsection{Transport-Based Topology}

Optimal-transport distances $W_{2}$ between persistence diagrams provide a
high-sensitivity probe of RDCC-induced non-Gaussian structure. These distances
capture global topological rearrangements that are difficult to mimic with
standard cosmological parameters.

% ========================================================
% 4 EUCLID-LEVEL FORECASTS
% ========================================================

\section{Euclid-Level Forecasts}

Using Euclid-like noise, resolution, and tomographic binning, RDCC forecasts
yield
\begin{equation}
    \sigma(\alpha) \sim \mathcal{O}(10^{-2}),
\end{equation}
with the strongest constraints arising from:

\begin{itemize}
    \item transport-based topological statistics,
    \item peak and void counts,
    \item neural-compressed summaries,
    \item combined multi-probe likelihoods.
\end{itemize}

Fisher-channel analysis shows that approximately $70\%$ of RDCC information
resides in the leading channel. This enables efficient compression and robust
inference even in high-dimensional summary-statistics spaces.

% ========================================================
% 5 DEGENERACY STRUCTURE AND EUCLID SYNERGY
% ========================================================

\section{Degeneracy Structure and Euclid Synergy}

RDCC variations can partially mimic $\Lambda$CDM parameter shifts. Euclid's
multi-probe design helps break these degeneracies:

\begin{itemize}
    \item \textbf{Tomographic binning} separates redshift-dependent relaxation effects.
    \item \textbf{Cross-correlation with galaxy clustering} reduces null directions.
    \item \textbf{Topological probes} provide strong orthogonality to
    $\Lambda$CDM-aligned manifolds.
\end{itemize}

The combination of geometric, morphological, and topological probes yields a
highly complementary constraint structure.

% ========================================================
% 6 INTEGRATION INTO EUCLID PIPELINES
% ========================================================

\section{Integration into Euclid Pipelines}

RDCC can be tested within Euclid's existing analysis framework through modular
extensions at the forward-modeling, summary-statistics, and likelihood levels.

\begin{itemize}
    \item \textbf{Forward modeling:}  
    RDCC can be incorporated into lensing simulations via modified potential
    kernels $\Phi_{\mathrm{RDCC}}(k,z)$.

    \item \textbf{Summary statistics:}  
    RDCC signatures are compatible with Euclid's power-spectrum pipeline and can
    be extended to include peaks, Minkowski functionals, persistent homology,
    and transport-based topology.

    \item \textbf{Likelihood integration:}  
    RDCC likelihood components can be added modularly to Euclid's cosmological
    parameter estimation framework.

    \item \textbf{Validation:}  
    The RDCC Benchmark Suite provides simulation-based calibration and
    cross-probe consistency tests.
\end{itemize}

% ========================================================
% 7 CONCLUSION
% ========================================================

\section{Conclusion}

RDCC provides a theoretically motivated and observationally testable extension of
$\Lambda$CDM. Its geometric and topological signatures are well matched to
Euclid's weak-lensing capabilities, and Euclid-level forecasts demonstrate that
RDCC can be constrained with percent-level precision. This paper outlines a
practical pathway for integrating RDCC into Euclid's analysis pipelines and
provides a survey-ready framework for testing relaxation-driven cosmology.

\end{document}